\label{chap:83-02}

\section{Réalisations géométriques du TRIT : régimes elliptique, parabolique et hyperbolique et double orientation temporelle}

Le TRIT conserve dans un même état le résidu publié, le cocycle de
retenue, l’orientation et la mémoire de composition. Ses \(3\) régimes
algébriques ne sont pas des couches éditoriales : ce sont des réalisations corrélatives de
l’opérateur quadratique \(J_{\tau}^{2}=-\tau I\), reliées aux \(2\)
orientations du temps et à l’état de seuil.

\begin{HMTScientificOwnerSubordinated}
\def\HMTTRITFormal{1}
\def\HMTTRITDevelopments{1}
\input{sections/hmt/02_trit_geometrias}
\let\HMTTRITFormal\undefined
\let\HMTTRITDevelopments\undefined
\end{HMTScientificOwnerSubordinated}

% RESULTADO_RECUPERADO. Owner integro: sections/hmt/09b_algebras_cuadraticas.tex
% Destino causal: capitulo 2; realizaciones cuadraticas y generadores TRIT.
\begin{HMTScientificOwnerSubordinated}
\input{sections/hmt/09b_algebras_cuadraticas}
\end{HMTScientificOwnerSubordinated}

% RESULTADO_RECUPERADO. Owner integro: sections/hmt/03k_genealogia_operaciones_rev2.tex
% Destino causal: capitulo 2; operaciones levantadas y memoria del TPK.
\begin{HMTScientificOwnerSubordinated}
\input{sections/hmt/03k_genealogia_operaciones_rev2}
\end{HMTScientificOwnerSubordinated}

\section[APP, TRIT et TPK : support, orientation et transport]
{APP, TRIT et TPK : support, orientation et transport}

\subsection{L’APP comme arithmétique des fibres}

L’APP n’est pas une table illustrative : c’est le support arithmétique où devient
visible la différence entre agréger et composer. Dans
\(
X_9=(\Znine)^2
\)
la somme accumule des déplacements et des histoires alternatives ; le produit compose
des étapes et, lorsqu’un facteur n’est pas inversible, replie des entrées distinctes sur
un même résidu. Le relèvement résidu--quotient déploie ces fibres et
conserve la généalogie contractée par la publication modulaire.

Soit \(q(a,b)=a+b\pmod 9\). La donnée publiée \(q(a,b)\) ne détermine pas la paire
\((a,b)\). La fibre
\[
 q^{-1}(r)=\{(a,b):a+b\equiv r\pmod9\}
\]
est la mémoire des alternatives. Pour la multiplication, la contraction est encore
plus marquée : les classes non unitaires \((0,3,6)\) produisent des fibres de tailles
différentes. L’APP publie donc une géométrie des opérations avant
toute interprétation physique.

Le bloc central actif est
\[
B_c=\begin{pmatrix}7&2\\2&7\end{pmatrix}
   =7I+2S=9P_++5P_-,
\qquad
P_\pm=\frac{I\pm S}{2}.
\]
Son saut spectral local est \(9-5=4\) ; le coefficient de couplage est \(2\).
La valeur \(54\) appartient à un déploiement ultérieur et, dans la fibre électronique,
à \(D_1(\Gamma_e)\). Les égaler parce qu’ils apparaissent dans une chaîne numérique effacerait
leurs types.

\subsection{Géométrie locale de Lorentz et plan de Minkowski \texorpdfstring{\(1+1\)}{1+1}}

La normalisation
\[
M_c=\frac{B_c}{\sqrt{45}}
=\exp\!\left(\frac12\log\frac95\,S\right)
\]
réalise un élément de \(SO^+(1,1)\). L’APP contient ainsi, avant la physique
continue, une décomposition en modes \(P_+,P_-\), un écart spectral, un couplage
et une rapidité. Le résultat n’est pas encore l’espace-temps de Minkowski : c’est
la structure algébrique qui rendra possible une représentation lorentzienne.

\subsection{TRIT comme orientation et mémoire locale}

TRIT introduit le signe que la publication modulaire omet. Dans une somme triadique
\[
 a+b=r+3c,
\]
le résidu \(r\) publie le résultat et la retenue \(c\) conserve la manière dont il a été
obtenu. Les états \((+1,0,-1)\) ne doivent pas être interprétés seulement comme \(3\)
couleurs. Dans la paramétrisation temporelle, ils agissent respectivement comme
retour avec mémoire du passé, seuil présent et ouverture bidirectionnelle vers le futur ;
dans la réalisation algébrique, ils distinguent contre-flux, fixation et flux ; dans la
réalisation géométrique, ils étiquettent les branches hyperbolique, parabolique et elliptique
d’une famille quadratique. La compatibilité entre ces lectures exige des applications,
non une identification verbale.

\begin{definition}[État TRIT enrichi]
Un état TRIT est une paire \((r,c)\) accompagnée d’une orientation. La première
composante est publiable ; la seconde appartient à la fibre de mémoire. Le
changement de lecteur peut oublier \(c\), mais le transport TPK ne peut pas le faire.
\end{definition}

Cette séparation permet de formaliser le présent, le passé et le futur sans les transformer
en trois instants absolus. Le présent est le lieu de couplage de deux
orientations ; le passé et le futur sont des sens de transport. L’inversion de
route échange les deux sens et laisse le point de couplage comme section
de frontière.

\subsection{TPK comme système de transport}

TPK réunit support, orientation et mémoire. Ses objets sont des sections préparées
finies ou coinductives ; ses morphismes sont des routes enregistrées compatibles avec les
lecteurs. Chaque route contient
\[
\Gamma=(\text{origine},\text{destination},\text{feuillet},\text{orientation},
\text{retenue},\text{frontière},\text{registre}).
\]
La composition concatène des chemins et compose des mémoires. L’inversion renverse
l’ordre et l’orientation. Les lecteurs, le calendrier de \(108\) phases, la frontière et le
registre sont des opérateurs internes ; non des observateurs externes ajoutés.

Une opération publiée \(f:X\to Y\) admet un relèvement
\[
 \widetilde f:X\longrightarrow Y\times M_f,
 \qquad \pi_Y\circ\widetilde f=f.
\]
Étant donnés des relèvements fixés de \(f\) et de \(g\), la composition a un
relèvement canonique associé
\[
 \widetilde{g\circ f}:X\longrightarrow Z\times M_f\times M_g.
\]
La mémoire n’est pas un commentaire de l’opération : elle fait partie de son codomaine
enrichi.

\begin{proposition}[Généalogie des opérations]
Dans TPK, somme, produit, puissance, racine, logarithme, dérivation et intégration
admettent une lecture commune :
\begin{center}
\begin{tabular}{@{}p{.18\textwidth}p{.31\textwidth}p{.38\textwidth}@{}}
\toprule
Opération & Action & Mémoire préservée \\
\midrule
Somme & agrégation des alternatives & provenance et retenue \\
Produit & composition des étapes & ordre et fibre repliée \\
Puissance & itération dans un monoïde & profondeur et orbite \\
Racine & inversion multivoque & branche et orientation \\
Logarithme & coordonnée d’itération & générateur et jauge \\
Dérivation & différence par échelle & bord et orientation \\
Intégration & composition des accroissements & chemin et conditions de frontière \\
\bottomrule
\end{tabular}
\end{center}
\end{proposition}

La puissance n’introduit pas une opération ontologiquement déconnectée de la somme et du
produit : elle itère la composition. La racine inverse cette itération et doit conserver
la sélection de branche. Le logarithme enregistre la profondeur. Cette hiérarchie explique
pourquoi les puissances nonadiques et leurs lacunes sont informatives pour l’APP :
l’opération élevée conserve des couches de mémoire qu’une racine numérique isolée
contracte.

\subsection{Périodicité temporelle de (108) phases et lois de retour à (27), (54) et (108) itérations}

Le calendrier observable possède trois seuils distincts. Après 27 itérations,
la position revient et l’orientation s’inverse ; après 54, la position
et l’orientation sont restaurées ; après 108, la phase observable complète est restaurée. La mémoire
du semi-groupe positif, cependant, n’est pas effacée. Schématiquement,
\[
F_{\rm obs}^{27}(x,o)=(x,-o),\qquad
F_{\rm obs}^{54}(x,o)=(x,o),\qquad
F_{\rm obs}^{108}(x,o,\phi)=(x,o,\phi).
\]
Cet ordre hiérarchique sera crucial pour séparer propagation à l’aller, retour,
orientation et cycle complet. Une période observable finie peut coexister avec
une mémoire généalogique non périodique.

\begin{falsadorBox}[title={Contrôles de type}]
Il n’est pas admissible d’identifier TPK à une liste de 34 opérateurs ; d’appeler
« catégorie fibrée » la structure sans relèvements cartésiens ; de situer la
retenue simultanément comme cobord cellulaire et comme 2-cocycle sans séparer leurs
complexes ; ou de transformer le bloc \(SO^+(1,1)\) en espace-temps physique avant
de construire la représentation métrique.
\end{falsadorBox}


\section{Fibration de Hopf, cercles de Villarceau, revêtement double et holonomie de Möbius}

\subsection{Structure complexe et fibration de Hopf}

Soit \(S^3\subset\mathbb C^2\) et
\[
h(z_0,z_1)=
\bigl(2\Re(z_0\overline z_1),
2\Im(z_0\overline z_1),|z_0|^2-|z_1|^2\bigr)
\in S^2.
\]
Dans le tore de Hopf
\[
T_\eta=\{(z_0,z_1):|z_0|=\cos\eta,
|z_1|=\sin\eta\},
\]
les courbes
\[
\gamma_+(t)=(\cos\eta\,e^{it},\sin\eta\,e^{it}),
\]
\[
\gamma_-(t)=(\cos\eta\,e^{it},\sin\eta\,e^{-it})
\]
présentent des comportements distincts : \(\gamma_+\) est une fibre de Hopf et
\(\gamma_-\) se projette avec un degré deux. Sous projection stéréographique, les
classes homologiques \((1,1)\) et \((1,-1)\) produisent les deux familles diagonales
de Villarceau.

La conjugaison complexe totale
\[
K(z_0,z_1)=(\overline z_0,\overline z_1)
\]
préserve chaque famille non orientée et inverse son orientation. Elle n’échange pas
les deux familles. La conjugaison partielle peut les échanger, mais ce n’est ni une
application complexe linéaire ni l’antiunitaire standard global. Cette correction
est essentielle pour ne pas fabriquer une symétrie particule--antiparticule à partir du
dessin toroïdal.

\subsection{Revêtement double et retour spinoriel}

L’application \(SU(2)\to SO(3)\) est un revêtement double. Un tour de
\(2\pi\) dans \(SO(3)\) peut se relever en \(-I\) dans \(SU(2)\) ; le tour de
\(4\pi\) récupère \(I\). La structure de signe ne transforme pas un cercle en
ruban de Möbius. Pour obtenir un ruban, il faut un fibré réel en droites
dont le caractère d’holonomie soit \(-1\).

Soit \(L_\chi\to S^1\) le fibré associé à
\[
\chi:\pi_1(S^1)\longrightarrow\{\pm1\},
\qquad \chi(1)=-1.
\]
Son espace total est un ruban de Möbius. La route \(C_M=-\operatorname{rev}\)
apporte le candidat de signe ; le cercle de Villarceau apporte le lacet ;
l’entrelaceur physique doit prouver que le premier est l’holonomie du second.

\subsection{Décomposition spectrale de l’opérateur orienté de mémoire et dérivation du rapport (5{:}2)}

Avec \(B_c=7I+2S\), définissons l’opérateur orienté de mémoire
\[
B_{\rm mem}^{\varepsilon}=5I+2\varepsilon S,
\qquad\varepsilon=\pm1.
\]
Pour \(\varepsilon=+1\),
\[
\Spec(B_{\rm mem}^{+})=\{7,3\},
\qquad\det B_{\rm mem}^{+}=21.
\]
Interpréter les valeurs \((7,3)\) comme les bords équatoriaux extérieur et intérieur
dans une carte positive commune produit
\[
R+r=7\ell,\qquad R-r=3\ell,
\]
et donc
\[
R=5\ell,\qquad r=2\ell,\qquad
\sin\beta=\frac rR=\frac25.
\]
L’angle de Villarceau est
\[
\beta=\arcsin\frac25
=23.5781784782\ldots^\circ.
\]

Le théorème algébrique \(\{7,3\}\leftrightarrow(5,2)\) est exact. La prémisse
géométrique qui interprète le spectre comme les deux bords d’un tore est une
interface déclarée. Elle ne doit pas être dissimulée et l’angle ne doit pas être employé comme cible.

La réalisation du bloc complet \(B_c\) conserve une autre lecture :
\(R:r=7:2\) et \(\beta=\arcsin(2/7)\). Ce ne sont pas des valeurs rivales ; elles relèvent de
deux foncteurs : bloc total et mode de mémoire. La priorité physique du second
doit être décidée par l’entrelaceur avec la dynamique, et non par une proximité décimale avec
une autre constante.

\subsection{Coordonnées angulaires conjuguées et normalisation de l’unité de phase dans la fibration de Hopf}

L’angle natif de HMT est une phase \(u\in\mathbb R/\mathbb Z\). Les lecteurs
publient
\[
\theta_{\rm rad}=2\pi u,
\qquad
\theta_{\rm deg}=360u.
\]
Radians et degrés sont des cartes distinctes du même élément circulaire. Le radian
n’est pas plus « ontologique » parce qu’il est adimensionnel ; sa naturalité provient de la
longueur d’arc. Les degrés fournissent une autre coordonnée absolue du
tour. Toute relation avec \(\alpha^{-1}\), \(\varphi\) ou CKM doit
d’abord être formulée comme un quotient ou une holonomie invariante, puis seulement
être évaluée dans une carte angulaire.

L’identité exacte déjà disponible
\[
\tanh^2(2\log\varphi)=\frac59
\]
relie le canal Perron--Fibonacci au quotient spectral de \(B_c\). Son
complément est \(4/9\), le saut normalisé. Cette relation est structurelle.
Elle n’implique pas que tout angle construit avec \(\varphi\) soit un angle physique.

\subsection{Brouwer, Möbius et degré}

Le théorème du point fixe de Brouwer, les transformations fractionnaires de
Möbius et le ruban de Möbius sont trois objets distincts. Ils ne sont reliés
qu’après fixation du domaine, de l’action et du fibré. Les orbites elliptiques d’avance et de
retard peuvent être modélisées comme deux sections d’un fibré à holonomie de
signe ; la transformation de Möbius décrit l’action conforme dans la carte ;
le théorème de Brouwer garantit un point fixe pour les applications continues du disque.
Aucune de ces phrases ne remplace le diagramme de réalisation.

\begin{fronteraBox}[title={Résultat et frontière}]
On démontre le degré double de la courbe de Hopf, l’orientation de la
conjugaison totale et la construction algébrique \(7,3\mapsto5,2\). La lecture
toroïdale (5{:}2), l’holonomie de Möbius et l’interprétation
particule--antiparticule forment une chaîne d’interfaces qui doit partager un
unique entrelaceur. Le volume ne force pas cet entrelaceur au moyen d’une
coïncidence angulaire.
\end{fronteraBox}


\section{Réalisation hilbertienne finie du système de chemins et de ses projecteurs conditionnels}

\subsection{Réalisation matricielle de chaque profondeur}

Pour \(d\ge1\), nous définissons
\[
  \mathcal C_9^{(d)}=(\mathbb Z/9\mathbb Z)^d,
  \qquad
  H_d=\ell^2(\mathcal C_9^{(d)}),
  \qquad
  A_d=B(H_d)\cong M_{9^d}(\mathbb C).
\]
La base canonique \(\{\delta_x:x\in\mathcal C_9^{(d)}\}\) conserve
l’étiquette discrète de chaque état. La mesure uniforme induit la trace
normalisée
\[
  \tau_d(a)=\frac{1}{9^d}\Tr(a).
\]

La factorisation
\[
  \mathcal C_9^{(d+1)}
  =
  \mathcal C_9^{(d)}\times\mathbb Z/9\mathbb Z
\]
produit une identification unitaire
\[
  H_{d+1}\cong H_d\otimes\mathbb C^9.
\]
C’est l’entrée exacte de la tour opératorielle. Cela n’identifie pas pour autant
\((\mathbb Z/9)^3\) à \(\mathbb F_3^6\) comme groupes : ils partagent \(729\)
éléments, mais leurs opérations et cocycles de retenue sont distincts.

\subsection{Représentation de Weyl–Heisenberg de la phase et du déplacement sur la fibre tritique}

Soit \(\omega=e^{2\pi i/9}\). Sur
\(H_1=\ell^2(\mathbb Z/9\mathbb Z)\), nous définissons
\[
  (Xf)(r)=f(r-1),
  \qquad
  (Zf)(r)=\omega^r f(r).
\]
Alors
\[
  X^9=Z^9=I,
  \qquad
  ZX=\omega XZ.
\]
Le groupe engendré par \(X,Z,\omega I\) est le groupe de Heisenberg fini
d’ordre \(9^3=729\).

\begin{theorem}[Stone--von Neumann fini]
Toute représentation unitaire irréductible du groupe de Heisenberg fini dans
laquelle le centre agit par le caractère primitif
\(\omega I\) est unitairement équivalente à la représentation précédente.
\end{theorem}

\begin{proof}
Soit \(v\) un vecteur propre de \(Z\). Les vecteurs
\[
  v,Xv,\ldots,X^8v
\]
ont des valeurs propres distinctes pour \(Z\), puisque
\(ZX^kv=\omega^kX^kZv\). Ils sont orthogonaux et engendrent un sous-espace
invariant. Par irréductibilité, ils forment une base de toute la représentation.
Dans cette base, \(X\) est le déplacement cyclique et \(Z\) la phase diagonale.
\end{proof}

\begin{controlbox}
L’unicité échoue si le caractère central primitif n’est pas fixé ou si l’on
permet des représentations réductibles. Ce théorème fini ne doit pas être confondu
avec la version continue pour les relations canoniques de commutation.
\end{controlbox}

\subsection{Projecteurs conditionnels et incompatibilité pré-Heisenberg}

Les feuillets additif et multiplicatif de \(\APP\) définissent des sous-espaces
d’information et leurs projecteurs orthogonaux
\(P_{\Sigma,d},P_{\Pi,d}\). Les identités construites satisfont
\[
 \|[P_{\Sigma,2},P_{\Pi,2}]\|
 =\frac{\sqrt2}{3},
 \qquad
 \|[P_{\Sigma,3},P_{\Pi,3}]\|
 =\frac{\sqrt3}{4}.
\]
La non-commutativité ne se déduit pas d’une analogie quantique : c’est un théorème
matriciel fini. Sa lecture « pré-Heisenberg » est postérieure à ce calcul.

\begin{proposition}
Les projecteurs précédents n’appartiennent pas conjointement à une même
sous-algèbre booléenne de projections.
\end{proposition}

\begin{proof}
Dans une sous-algèbre abélienne, tous les projecteurs commutent. Les commutateurs
montrés sont non nuls.
\end{proof}

\subsection{Structure hilbertienne de la publication matricielle et mémoire généalogique préservée}

L’espace de Hilbert fournit la superposition linéaire, la projection orthogonale
et le calcul spectral. HMT apporte l’indice généalogique de la base et distingue
quel opérateur représente une transition, une mesure ou une ombre. La
réalisation ne remplace pas le \(\TPK\) : elle le représente.

\begin{sourcebox}
Les projecteurs APP et la paire Weyl--Heisenberg font partie de la construction
précédente. L’organisation \(H_d,A_d\) et la preuve d’unicité sont établies
ici par l’application du théorème classique correspondant aux objets
HMT.
\end{sourcebox}


\section{Structure symplectique finie de l’APP et fibration nonadique des observables}
\label{p12-83-c2-s4:gmc:chap:simplectica}

La géométrie symplectique ne fait pas partie des théorèmes publiés sur
les carrés magiques de De las Cuevas--Netzer. Le lien légitime passe par une
voie voisine et standard : opérateurs de Weyl, espace des phases discret et
bases mutuellement non biaisées \citep{gibbons2004phase,planat2007rings}. Dans HMT,
cette voie acquiert un relèvement nonadique et une mémoire de fibre.

\subsection{La forme alternée et la commutation de Weyl}

Sur $V_3=\mathbb F_3^2$, nous définissons
\[
 \langle(q,p),(q',p')\rangle_{\rm sp}=qp'-pq'.
 \tag{GMC-5.1}
\]
Si $X|k\rangle=|k+1\rangle$,
$Z|k\rangle=\omega^k|k\rangle$ et $\omega=e^{2\pi i/3}$, les opérateurs
$W_{q,p}=X^qZ^p$ satisfont, avec ces conventions,
\[
 W_uW_v=\omega^{-\langle u,v\rangle_{\rm sp}}W_vW_u.
 \tag{GMC-5.2}
\]
Par conséquent, deux opérateurs commutent exactement lorsque leurs étiquettes sont
orthogonales pour la forme alternée. En dimension deux, les sous-droites
isotropes maximales sont les quatre points de
\[
 \mathbb P^1(\mathbb F_3)=\{[1:0],[0:1],[1:1],[1:2]\}.
 \tag{GMC-5.3}
\]
Chacune engendre un contexte abélien maximal de $M_3(\mathbb C)$ ; ses
projecteurs spectraux constituent une \PVM. Les contextes associés à des droites
distinctes produisent les quatre bases MUB de \textup{(GMC-2.13)}.

\begin{proposition}[Décomposition qutrit par directions]
\label{p12-83-c2-s4:gmc:prop:descomposicion-masa-qutrit}
Soient $\mathcal D_a$ les quatre sous-algèbres diagonales déterminées par les
droites de \textup{(GMC-5.3)} et $\mathcal D_a^0$ leurs parties de trace nulle. Alors
\[
 M_3(\mathbb C)
 =\mathbb CI\oplus
   \bigoplus_{a\in\mathbb P^1(\mathbb F_3)}\mathcal D_a^0
 \tag{GMC-5.4}
\]
comme somme orthogonale pour le produit de Hilbert--Schmidt. Le compte des
dimensions est $9=1+4\cdot2$.
\end{proposition}

\begin{proof}
Chaque $\mathcal D_a^0$ est de dimension deux. L’absence de biais implique
$\Tr(A^*B)=0$ pour $A\in\mathcal D_a^0$ et
$B\in\mathcal D_b^0$, $a\ne b$. Toutes sont orthogonales à $I$. Les
dimensions totalisent neuf, la dimension de $M_3(\mathbb C)$.
\end{proof}

Cette identité explique pourquoi les douze projections ne sont pas douze pièces
indépendantes arbitraires : ce sont quatre résolutions de l’unité dont les parties
centrées décomposent tout l’espace opératoriel qutrit.

\subsection{Le relèvement \texorpdfstring{$12\to4\times3$}{12 vers 4 fois 3}}

Sur le module nonadique $V_9=(\mathbb Z/9\mathbb Z)^2$, la forme
alternée est conservée
\[
 \langle(a,b),(c,d)\rangle_9=ad-bc\pmod9.
 \tag{GMC-5.5}
\]
La droite projective $\mathbb P^1(\mathbb Z/9\mathbb Z)$ est formée des
vecteurs primitifs, modulo la multiplication par une unité. Elle peut
être représentée comme
\[
 \{[1:t]:t\in\mathbb Z/9\mathbb Z\}
 \ \sqcup\ 
 \{[3u:1]:u\in\mathbb Z/3\mathbb Z\},
 \tag{GMC-5.6}
\]
et possède $9+3=12$ points. La réduction modulo trois donne le diagramme
\[
 \mathbb P^1(\mathbb Z/9\mathbb Z)
 \xrightarrow{\ r\ }
 \mathbb P^1(\mathbb F_3),
 \qquad |r^{-1}(a)|=3.
 \tag{GMC-5.7}
\]

La base $a$ d’un contexte qutrit et le résultat $k$ de sa \PVM
forment également un ensemble $4\times3$. Un \emph{étalonnage} est une bijection
\[
 \Xi:\mathbb P^1(\mathbb Z/9\mathbb Z)
 \longrightarrow
 \{P_{a,k}:a\in\mathbb P^1(\mathbb F_3),\ k\in\mathbb F_3\}
 \tag{GMC-5.8}
\]
qui envoie chaque fibre de $r$ vers les trois résultats du contexte situé au-dessus de
$a$.

\begin{controlbox}
La fibration est canonique ; l’affectation de l’ordre interne de chaque fibre aux
trois résultats $k$ est une jauge. Une bijection n’est pas encore une
équivariance. Pour promouvoir $\Xi$ en entrelaceur, il faut démontrer que les
actions linéaires, projectives et antiunitaires pertinentes sont transportées dans
toutes les fibres. La jauge préliminaire échoue à une vérification antiunitaire dans une
fibre ; cet échec est conservé comme réfutateur, il n’est pas dissimulé.
\end{controlbox}

\subsection{Profondeur, retenue et mémoire sur l’espace des phases discret}

La géométrie de Weyl sur $\mathbb F_3^2$ explique la commutation, les contextes et les
MUB. HMT ajoute trois niveaux qui ne peuvent être retrouvés à partir de la seule forme
symplectique :

\begin{enumerate}
 \item le relèvement d’une direction ternaire vers trois points de
 $\mathbb P^1(\mathbb Z/9\mathbb Z)$ ;
 \item le feuillet, la retenue et l’orientation qui distinguent les états de
 même réduction ;
 \item la route TPK qui transporte ces coordonnées et peut revenir à la
 phase sans revenir à l’état.
\end{enumerate}

L’apport potentiel à la littérature qutrit n’est pas « une autre construction de
quatre MUB », déjà connue. C’est une géométrie de \emph{racines de contexte} : chaque
contexte ternaire possède trois antécédents nonadiques, et ces antécédents
peuvent porter une dynamique et une monodromie. Le carré d’ordre douze du
\Cref{p12-83-c22-s3:gmc:thm:qms-12-3} est la première ombre normalisée de cet atlas.

\subsection{Séparation typée entre le recouvrement qutrit
\texorpdfstring{$1/3$}{1/3} et le lecteur HMT
\texorpdfstring{$1/\pi$}{1/pi}}

Dans une paire de bases MUB qutrit,
\[
 \Tr(P_{a,k}P_{b,\ell})=\frac13.
 \tag{GMC-5.9}
\]
Dans le lecteur HMT défini ici, la relation entre un feuillet réel et une
lecture volumétrique peut prendre le quotient
\[
 \frac{\pi}{\pi^2}=\frac1\pi.
 \tag{GMC-5.10}
\]
Les deux expressions ne sont pas interchangeables : $1/3$ est un recouvrement entre
projecteurs en dimension trois ; $1/\pi$ est un poids scalaire d’un autre lecteur.
La différence n’empêche pas de les coordonner. Une \PVM
$\{P_0,P_1,P_2\}$ étant fixée, nous définissons une mesure à neuf effets :
\[
 F_k^{(\pi)}=\frac3\pi P_k\quad(k=0,1,2),
 \qquad
 F_r^{(\pi)}=\frac{1-3/\pi}{6}I_3\quad(r=3,\ldots,8).
 \tag{GMC-5.11}
\]
C’est une \POVM parce que $3/\pi<1$ et
\[
 \sum_{r=0}^{8}F_r^{(\pi)}
 =\frac3\pi I_3+\left(1-\frac3\pi\right)I_3=I_3.
 \tag{GMC-5.12}
\]
Si $Q$ est un projecteur de rang un d’un contexte MUB distinct,
\[
 \Tr(QF_k^{(\pi)})
 =\frac3\pi\Tr(QP_k)
 =\frac3\pi\cdot\frac13
 =\frac1\pi.
 \tag{GMC-5.13}
\]

\begin{bridgebox}
La formule \textup{(GMC-5.11)} transforme exactement la probabilité MUB $1/3$ en
le lecteur $1/\pi$ au moyen de la jauge $3/\pi$, sans identifier les deux nombres.
Les six effets résiduels conservent l’unité et complètent une \POVM à
neuf résultats, à partir de laquelle on obtient un autre \QMS cyclique $(9,3)$. Cette
construction n’identifie pas encore la provenance géométrique de la masse
résiduelle ; elle fixe le type opératoriel exact qu’elle devra avoir.
\end{bridgebox}

\subsection{Groupe symplectique et équation d’équivariance}

En dimension deux sur $\mathbb F_3$, le groupe symplectique coïncide avec
$SL(2,\mathbb F_3)$ et permute les quatre droites. Son action se relève,
projectivement, par le groupe de Clifford sur les \PVM qutrit. Soit
$G_{\rm HMT}$ un groupe ou groupoïde de transformations de l’atlas nonadique et
$U_g$ l’opérateur unitaire ou antiunitaire candidat dans la carte qutrit. La
condition qui transforme un étalonnage en un pont dynamique est
\[
 \Xi(g\cdot\ell)=U_g\,\Xi(\ell)\,U_g^*,
 \qquad
 g\in G_{\rm HMT},\quad
 \ell\in\mathbb P^1(\mathbb Z/9\mathbb Z).
\tag{GMC-5.14}
\]
Si $g$ est une flèche et non seulement une symétrie globale, l'équation doit inclure
la source et le but et se formule comme une transformation naturelle. La vérification de
\textup{(GMC-5.14)} est l'étape qui distingue un étiquetage heureux d'une
transduction structurelle.

\begin{sourcebox}
Les identités symplectiques et MUB s'appuient sur la littérature relative à la phase finie
\citep{gibbons2004phase,planat2007rings}. Elles ne sont pas attribuées aux articles de
De las Cuevas--Netzer sur les carrés magiques. La fibration nonadique, l'étalonnage
et les lecteurs de route constituent la contribution propre de la
construction HMT étudiée ici.
\end{sourcebox}
